\begin{textsample}{POS Dim 3 – 2013\_c6 – Score -2.00 – 2013\_c6/enade\_2013\_med\_9c8209e0a3.txt}  \label{ex:f3_pos_231}
CAPÍTULO 3 ANÁLISE TÉCNICA DA PROVA Este capítulo tem por objetivo apresentar o desempenho dos estudantes concluintes de Medicina no ENADE/2013.
Para isso, foram calculadas as estatísticas básicas da prova em seu todo ( seção 3.1.1 ), bem como as estatísticas dos componentes relacionadas ao Componente de Formação Geral ( seção 3.1.2 ) e ao de Conhecimento Específico da Área ( seção 3.1.3 ).
Dadas as suas características, foram analisadas em separado, as questões objetivas ( seção 3.2 ) e as questões discursivas ( seção 3.3 ).
Para as questões objetivas, foram disponibilizados os índices de facilidade e o de discriminação bisserial, também em separado para os Componentes de Formação Geral ( seção 3.2.1 ) e de Conhecimento Específico ( seção 3.2.2 ).
De cada componente, uma das questões foi escolhida para exemplificar a análise gráfica relacionando as alternativas escolhidas pelos estudantes ( inclusive o gabarito ) com o número de acertos no componente.
O Anexo I apresenta a íntegra da análise gráfica para todas as questões objetivas.
Para cada uma das questões discursivas, os conteúdos dos tipos mais comuns de respostas dos alunos são apresentados e comparados com o padrão de resposta esperado ( ver Anexo VIII com o padrão de respostas ).
Tomando como base as duas questões discursivas do Componente de Formação Geral, a seção 3.3.1.6 apresenta comentários sobre a correção das respostas com respeito à Língua Portuguesa.
Nas tabelas, são apresentados o tamanho da população inscrita e de presentes, e as seguintes estatísticas das notas¹¹ : média do desempenho na prova, erro padrão da média, desvio padrão, nota mínima, mediana e nota máxima.
Tais estatísticas contemplam o total de estudantes concluintes da Área de Medicina inscritos e presentes à prova do ENADE/2013, tendo -se em vista as seguintes agregações : ( a ) as Grandes Regiões e o país como um todo ; ( b ) a Categoria Administrativa ; e ( c ) a Organização Acadêmica.
Em relação aos gráficos de distribuição de notas, o intervalo considerado foi de 10 unidades, aberto à esquerda e fechado à direita, com exceção do primeiro intervalo, [ 0 ; 10 ], fechado em ambos os extremos.
Para os gráficos de distribuição das notas das questões discursivas, foram consideradas mais duas categorias : questão em branco e nota zero. ¹¹ Essas estatísticas e outras estão definidas no Capítulo 1. 3.1 ESTATÍSTICAS BÁSICAS DA PROVA 3.1.1 Estatísticas Básicas Gerais A Tabela 3.1 apresenta as estatísticas básicas da prova por grande Região.
A população total de inscritos foi de 16.264.
Destes, 15.617 estiveram presentes, sendo 4,0 \% o índice de não comparecimento.
A Região de maior abstenção foi a Norte ( 19,6 \% ) e a de menor abstenção foi a Centro-Oeste ( 0,9 \% ).
A média das notas da prova como um todo ( nas seções seguintes serão analisados os Componentes de Formação Geral e de Conhecimento Específico ) foi 47,7, sendo que os alunos da região Norte obtiveram a média mais baixa ( 41,2 ) e os da região Centro-Oeste obtiveram a média mais alta ( 51,1 ).
As demais médias foram : 48,6 na região Nordeste, 47,4 na região Sudeste e 49,6 na região Sul.
O desvio padrão para o Brasil como um todo foi 13,4, sendo o maior desvio padrão encontrado na região Sul ( 15,4 ) e o menor nas regiões Norte e Nordeste ( 12,3 ), indicando uma menor dispersão das notas destas últimas regiões.
A região que obteve a maior nota máxima foi a Sul ( 89,6 ), ao passo que a região que atingiu a menor nota máxima foi a Norte ( 77,9 ).
A mediana do Brasil como um todo foi 47,7, sendo a maior mediana obtida nas regiões Sul e Centro-Oeste ( 50,8 ) e a menor obtida na Norte ( 41,3 ).
A nota mínima foi zero em todas as regiões, sem exceção.
O comportamento das notas dos estudantes de todo o Brasil pode ser observado no Gráfico 3.1 que apresenta um histograma com a distribuição das mesmas.
Essa é uma distribuição unimodal com moda no intervalo ( 40 ; 50 ].
O coeficiente de assimetria da distribuição das notas é negativo ( -0,22 ) e não podemos considerar que a distribuição é aproximadamente simétrica.
As distribuições para as regiões Norte, Nordeste e Centro-Oeste apresentam assimetria positiva, concentração pouco maior do lado esquerdo do histograma e mais espalhada do lado direito, enquanto que para a região Sul a distribuição apresenta assimetria negativa.
A região Sudeste apresenta um coeficiente de assimetria igual a zero.
Os Gráficos 3.2, 3.3 e 3.4 apresentam informações referentes à média da nota final dos participantes, desagregando os resultados de acordo com, respectivamente, as Grandes Regiões do país, a Categoria Administrativa e a Organização Acadêmica.
Os gráficos apresentam o valor da média das notas como uma barra e os extremos do intervalo de confiança de 95 \% como linhas verticais unidas por uma linha horizontal na forma da letra H maiúscula.
Considerando -se o gráfico de notas segundo Grande Região ( Gráfico 3.2 ), observa -se que existe diferença estatisticamente significativa ao nível de 95 \% entre a maior média, obtida na região Centro-Oeste ( 51,1 ) e a menor, obtida na região Norte ( 41,2 ).
Levando -se em conta os estudantes por Categorias Administrativas ( Gráfico 3.3 ), observa -se que existe diferença estatisticamente significante entre as médias das notas das IES Públicas e Privadas.
Pode ser observado ainda que a média dos alunos de IES Públicas ( 51,9 ) é maior que a dos alunos de IES Privadas ( 45,5 ).
A diferença entre as médias das regiões Centro-Oeste e Norte ( 9,9 ), a maior e a menor média, é superior a diferença entre IES Públicas e Privadas ( 6,4 ).
Tendo como foco o Gráfico 3.4, que apresenta as notas médias das provas segundo Organização Acadêmica, constata -se que a diferença da média das notas dos estudantes provenientes de Universidades é estatisticamente significativa ao nível de 95 \% das médias das notas dos estudantes provenientes de Centros Universitários e Faculdades, assim como a diferença entre a média dos estudantes de Centros Universitários e Faculdades.
A maior média foi obtida pelos estudantes de Universidades ( 48,3 ), e a menor, pelos de Centros Universitários ( 44,7 ). 3.1.2 Estatísticas Básicas no Componente de Formação Geral A Tabela 3.2 apresenta as estatísticas básicas em relação ao componente da prova que avalia a Formação Geral dos estudantes concluintes.
Os alunos de todo Brasil obtiveram desempenho médio de 56,0.
Quanto à variabilidade, o desvio padrão das notas dos estudantes do Brasil como um todo foi 16,2.
A maior média foi obtida na região Centro-Oeste ( 60,4 ), e a menor, na região Norte ( 51,8 ).
As demais médias foram : 57,2 na região Nordeste, 55,9 na região Sudeste e 55,4 na região Sul.
Já o maior desvio padrão foi obtido na região Sul ( 18,5 ) e o menor na região Centro-Oeste ( 13,9 ).
Os demais desvios padrões foram : 17,1 na região Norte, 15,7 na região Nordeste e 15,5 na região Sudeste.
A maior nota no Componente de Formação Geral da prova do ENADE foi obtida por pelo menos um aluno da região Sudeste ( 97,4 ) enquanto que a menor nota máxima foi obtida na região Centro-Oeste ( 91,0 ).
Nas outras regiões as notas máximas foram : 93,0 para a região Norte, 96,2 para região Nordeste e 94,8 para a região Sul.
A mediana do Brasil como um todo foi 58,0, sendo a menor mediana encontrada na região Norte ( 54,3 ) e a maior encontrada na região Centro-Oeste ( 61,6 ).
A nota mínima nesta parte foi zero em todas as regiões, sem exceção.
O Gráfico 3.5 propicia a avaliação do desempenho dos estudantes no Componente de Formação Geral a partir do histograma da distribuição das notas correspondentes.
A distribuição é unimodal, com moda nos intervalos ( 50 ; 60 ] e ( 60 ; 70 ], enquanto na prova como um todo a moda foi alcançada no intervalo ( 40 ; 50 ].
Nota -se, ainda, que no Gráfico 3.5 as notas apresentam uma maior dispersão do que no Gráfico 3.1 ( distribuição das notas da prova ), confirmado pela comparação dos desvios padrões : 13,4 para a nota da prova como um todo e 16,2 para o Componente de Formação Geral.
Para o Componente de Formação Geral, o coeficiente de assimetria da distribuição das notas dos estudantes é negativo ( -0,73 ), apresentando uma leve concentração à direita e cauda maior à esquerda.
Em todas as Grandes Regiões os histogramas também possuem assimetria negativa ( entre -1,21 para o Sul e -0,48 para o Norte ).
Nos Gráficos 3.6, 3.7 e 3.8 são apresentadas as informações referentes ao desempenho dos concluintes no Componente de Formação Geral, em diferentes agregações : Grande Região do país, Categoria Administrativa e Organização Acadêmica.
Observa -se pelo Gráfico 3.6 que existe diferença estatisticamente significativa ao nível de 95 \% entre a maior média das notas na Componente de Formação Geral, obtida na região Centro-Oeste ( 60,4 ) e a menor, obtida na região Norte ( 51,8 ).
Vemos que o intervalo de confiança mais largo é o da região Norte ; já o intervalo mais estreito é observado na região Sudeste.
Este fato está relacionado, também, com o tamanho da população envolvida, menor na região Norte do que na Sudeste.
No Gráfico 3.7, que representa as notas médias no Componente de Formação Geral segundo Categoria Administrativa do país, observa -se que existe diferença estatisticamente significativa entre as médias.
Assim como ocorreu para as notas da prova como um todo, para o Componente de Formação Geral, os concluintes de Medicina, das IES Públicas ( 58,7 ) obtiveram uma média maior do que os das IES Privadas ( 54,6 ).
Considerando -se o tipo de Organização Acadêmica, nota -se, no Gráfico 3.8, que existe uma diferença estatisticamente significativa entre a maior e a menor média.
Nas Universidades, a média ( 56,4 ) é mais elevada do que em Centros Universitários ( 54,0 ) e Faculdades ( 55,7 ). 3.1.3 Estatísticas Básicas do Componente de Conhecimento Específico A Tabela 3.3 apresenta as estatísticas básicas referentes ao Componente de Conhecimento Específico da área de Medicina.
A média do desempenho dos alunos do Brasil como um todo foi 44,9.
A maior média foi obtida na região Centro-Oeste ( 48,0 ), e a menor, na região Norte ( 37,7 ).
As demais médias foram : 45,7 na região Nordeste, 44,5 na região Sudeste e 47,6 na região Sul.
Quanto à variabilidade das notas, o desvio padrão do Brasil como um todo foi 15,0, sendo o maior desvio padrão observado na região Sul ( 16,5 ) e o menor na região Norte ( 13,1 ).
Os demais desvios padrões foram : 14,1 na região Nordeste, 14,8 na região Sudeste e 15,4 na região Centro-Oeste.
A mediana das notas dos estudantes de todo o Brasil foi 44,3.
A maior mediana ocorreu na região Sul ( 48,2 ) e a menor na região Norte ( 37,0 ).
As demais medianas foram : 44,6 na região Nordeste, 43,8 na região Sudeste e 46,8 na região Centro-Oeste.
A nota máxima do Brasil como um todo foi 92,8, sendo obtida por pelo menos um aluno da região Sul.
As demais notas máximas foram : 81,6 na região Norte, 88,2 na região Nordeste, 92,4 na região Sudeste e 88,8 na região Centro-Oeste.
A nota mínima foi zero em todas as regiões, sem exceção.
Assim como os Gráficos 3.1 e 3.5, o Gráfico 3.9, apresentado a seguir, proporciona uma avaliação do desempenho de concluintes em relação ao Componente de Conhecimento Específico com um histograma da distribuição das notas correspondentes.
Esta também é uma distribuição unimodal, e o grupo modal é o ( 40 ; 50 ], abaixo do grupo modal para a Formação Geral.
O coeficiente de assimetria da distribuição das notas do Componente de Conhecimento Específico é positivo ( 0,06 ), porém de pouca magnitude.
Nota -se pelo histograma ( Gráfico 3.9 ) que esta distribuição apresenta uma leve concentração à esquerda e cauda maior à direita.
Em quase todas as Grandes Regiões os histogramas também possuem assimetria positiva ( entre 0,31 para o Norte e 0,17 para o Sudeste ), exceto na região Sul ( -0,66 ) que possui assimetria negativa.
Os Gráficos 3.10, 3.11 e 3.12 apresentam, respectivamente, uma comparação dos resultados em relação à Grande Região do país, à Categoria Administrativa e à Organização Acadêmica, agora levando em conta o desempenho de alunos no Componente de Conhecimento Específico da prova.
Pelo Gráfico 3.10, observa -se que existe diferença estatisticamente significativa entre a maior e a menor média das notas no Componente de Conhecimento Específico, nas regiões Centro-Oeste ( 48,0 ) e Norte ( 37,7 ) respectivamente.
Quanto à Categoria Administrativa ( Gráfico 3.11 ), observa -se um comportamento semelhante àquele da Parte de Formação Geral e à prova como um todo, ou seja, existe diferença estatisticamente significativa entre as médias das IES Públicas e Privadas, sendo que a maior média foi obtida por alunos de IES Públicas de ensino ( 49,6 ).
Quanto ao Gráfico 3.12, observa -se, mais uma vez, que existe diferença estatisticamente significativa ao nível de 95 \% entre as notas no Componente de Conhecimento Específico de Universidades e dos demais tipos de Organização Acadêmica.
Sendo que a média dos concluintes das Universidades ( 45,6 ) foi maior do que de Centros Universitários ( 41,6 ) e de Faculdades ( 44,2 ). 3.2 ANÁLISE DAS QUESTÕES OBJETIVAS 3.2.1 Componente de Formação Geral A Tabela 3.4 apresenta as estatísticas básicas relativas às oito questões objetivas do componente da prova que abrange a Formação Geral dos estudantes.
A média do Brasil foi 60,1.
A menor média foi encontrada na região Norte ( 56,1 ) e a maior na região Centro-Oeste ( 64,0 ).
As demais médias foram 62,8 na região Nordeste, 59,2 na região Sudeste e 59,9 na região Sul.
O desvio padrão do Brasil foi 19,4, sendo o maior desvio padrão encontrado na região Sul ( 21,5 ) e o menor na região Centro-Oeste ( 17,1 ).
Os demais desvios foram : 20,0 na região Norte, 18,5 na região Nordeste e 19,0 na região Sudeste.
As mesmas medianas ( 62,5 ), notas máximas ( 100,0 ) e notas mínimas ( 0,0 ), nas questões objetivas do Componente de Formação Geral, foram obtidas em todas as regiões, sem exceção.
A Tabela 3.5 apresenta o índice de facilidade e o índice de discriminação ( ponto bisserial ) para cada uma das questões objetivas do Componente de Formação Geral.
Quanto ao índice de facilidade, foram usadas as seguintes cores para diferenciar o nível de dificuldade da questão :  Azul para as questões classificadas com índice muito fácil ( índice \&gt ; =0,86 ), verde para as questões classificadas com índice fácil ( 0,61 a 0,85 ), amarelo para as questões classificadas com médio ( 0,41 a 0,60 ), vermelho para as questões classificadas com difícil ( 0,16 a 0,40 ) e roxo para as questões classificadas com muito difícil ( \&lt ; =0,15 ).
Já quanto ao índice de discriminação, foram usadas as seguintes cores para qualificar a questão :  As questões classificadas com índice fraco receberam a cor vermelha ( índice \&lt ; =0,19 ), as classificadas com médio receberam a cor amarelo ( 0,20 a 0,29 ), as classificadas com bom receberam a cor verde ( 0,30 a 0,39 ) e as classificadas com muito bom ( \&gt ; =0,40 ) receberam a cor azul.
As questões objetivas do Componente de Formação Geral, segundo o índice de facilidade, foram assim avaliadas : das oito questões, uma teve o índice de facilidade classificado como muito fácil e duas como fácil.
Três questões foram tidas como médio, por terem índice de acertos situado na faixa entre 0,41 e 0,60 ( de 41,0 \% a 60,0 \% de acertos ).
Duas questões foram consideradas de dificuldade difícil, situando -se no intervalo entre 0,16 e 0,40 do índice de facilidade, ou seja, houve entre 16,0 \% e 40,0 \% de acertos.
Como já comentado, para análise das questões objetivas relativas à Formação Geral segundo o poder de discriminação, utilizou -se, o índice de discriminação ponto bisserial.
Nesta análise, as questões foram assim avaliadas : seis das oito questões apresentaram índices acima de 0,40 e, assim, foram classificadas com índice muito bom para esse grupo de alunos ; duas questões tiveram índice de discriminação bom, entre 0,30 e 0,39.
Nenhuma questão foi classificada com um índice fraco e, portanto, nenhuma questão foi eliminada pelo critério ponto bisserial.
O índice de facilidade variou de 0,29 a 0,96, e o de discriminação, de 0,35 a 0,49.
As questões com índice de discriminação muito bom, de números 3, 4, 5, 6, 7 e 8, ficaram espalhadas entre os diferentes níveis de facilidade : duas delas ( questões 4 e 8 ) classificadas na categoria fácil e três delas ( questões 3, 6 e 7 ) na categoria médio do índice de facilidade e uma questão ( 5 ) na difícil.
Em particular, a questão 7 apresentou maior poder discriminatório, com índice 0,49, e foi a quinta mais fácil, com uma proporção de 0,59 de acertos.
A questão de número 2 apresentou índice de facilidade 0,29, ou seja, um quantitativo de 29,0 \% de estudantes conseguiu resolvê -la, dentro do universo de participantes.
Além disso, seu índice de discriminação foi bom, com valor igual a 0,35.
O Gráfico 3.13, para exemplificar, analisa o comportamento da questão de número 3 de Formação Geral.
Trata -se da quarta questão mais fácil e a que obteve o segundo maior índice de discriminação dessa parte da prova.
Neste gráfico, cada uma das cinco curvas representa o percentual de respostas em determinada alternativa da questão, em função da nota dos estudantes nesta parte da prova ( Formação Geral/Múltipla Escolha ), antes de possíveis eliminações pelo critério do ponto bisserial.
A curva em vermelho corresponde à alternativa E, a correta para esta questão.
Assim, observa -se que entre os estudantes com notas mais baixas, nessa parte do exame, a situação mais frequente foi a escolha de uma das alternativas incorretas : a alternativa B ( em verde ) ou C ( em preto ).
Na medida em que a nota aumenta, indicando desempenho melhor nesta parte da prova, aumenta concomitantemente a proporção de estudantes que selecionaram a alternativa correta E, atingindo 100 \% para os alunos com 8 acertos.
Essa análise permite verificar como a questão discriminou os grupos de desempenho, justificando o alto índice obtido na questão.
Os gráficos relativos às demais questões de Formação Geral constam do Anexo I. 3.2.2 Componente de Conhecimento Específico A Tabela 3.6 apresenta as estatísticas básicas em relação às questões objetivas do Componente de Conhecimento Específico da prova, por Grande Região.
A média do Brasil deste componente foi de 48,8.
A menor média foi observada na região Norte ( 41,5 ) e a maior na região Centro-Oeste ( 51,7 ).
O desvio padrão de todo o Brasil foi 15,9, sendo o menor desvio padrão encontrado na região Norte ( 14,1 ) e o maior na região Sul ( 17,5 ).
A mediana de todo o Brasil foi 48,0.
Nas regiões Sudeste e Nordeste, a mediana foi também 48,0.
Na região Sul e na Centro-Oeste, a mediana foi 52,0, enquanto na região Norte a mediana foi 40,0.
A maior nota máxima da prova foi obtida, nas questões objetivas do Componente de Conhecimento Específico, por pelo menos um aluno tanto na região Sul quanto na Sudeste ( 100,0 ).
As demais notas máximas foram : 88,0 na região Norte, 96,0 na região Nordeste e 92,0 na região Centro-Oeste.
Em todas as regiões, sem exceção, a nota mínima foi zero.
A Tabela 3.7 apresenta o índice de facilidade e o índice de discriminação ( ponto bisserial ) para cada uma das questões objetivas do Componente de Conhecimento Específico.
Para facilitar a diferenciação das questões usaram -se as mesmas cores da Tabela 3.5 para as diferentes classificações dos índices de facilidade e de discriminação.
Dentre as questões objetivas da parte da prova relativa ao Componente de Conhecimento Específico, nenhuma foi anulada pela Comissão.
Desse modo, a classificação quanto ao índice de facilidade foi estabelecida com base nas 27 questões.
A partir dos índices obtidos, pode -se concluir que grande parte das 27 questões objetivas da prova foi considerada como difícil, 11 ao todo.
As 16 questões restantes foram classificadas com índice de facilidade da seguinte maneira : oito fáceis, sete médias e uma muito difícil.
Já quanto aos índices de discriminação das questões objetivas do Componente de Conhecimento Específico da prova, tem -se como resultado a seguinte classificação : dez das 27 questões foram consideradas como boas, enquanto seis delas tiveram índice de discriminação muito bom.
Assim, para 16 das 27 questões, os índices de discriminação foram bons ou muito bons.
Dentre as demais, nove delas foram classificadas como médias e outras duas como fracas, sendo 11, por conseguinte, a quantidade de questões nos dois patamares mais baixos de discriminação.
Constata -se, assim, que a prova – no que se refere ao Componente de Conhecimento Específico – possuía capacidade boa de discriminar entre aqueles que dominam ou não o conteúdo.
Dentre as questões que alcançaram os maiores índices de discriminação, as de números 11, 12, 17, 18, 25 e 35 classificadas com índice muito bom, igual ou superior a 0,40, duas delas ( questões 11 e 25 ) foram classificadas na categoria fácil, duas ( questões 18 e 35 ) como média e outras duas como difícil, quanto ao índice de facilidade.
A questão de número 24 foi a mais difícil dentre as 27 questões específicas, com baixo índice de facilidade, apenas 5,0 \% de acertos.
Essa questão apresentou poder discriminatório igualmente baixo, 0,06, o que comprova ter sido esta a mais difícil para os estudantes.
Destaca -se, também, a questão 16, com índice de facilidade 0,37, o que, em termos percentuais, corresponde a 37,0 \% de estudantes que responderam acertadamente, obtendo, ainda, 0,17 de índice de discriminação.
As questões 16 e 24 foram as únicas eliminadas do cômputo da nota final pelo critério do ponto bisserial.
A título de exemplo das análises do comportamento das questões objetivas, o Gráfico 3.14 analisa a questão 35 do Componente de Conhecimento Específico.
Esta foi uma das questões classificadas como de dificuldade Média nesta prova, apresentando índice de facilidade 0,51, ou seja, 51,0 \% dos estudantes assinalaram acertadamente a Seu índice de discriminação foi igual a 0,48, classificado como muito bom.
Neste gráfico, cada uma das cinco curvas representa o percentual de respostas em determinada alternativa da questão 35, em função da nota dos estudantes nesta parte da prova, antes de possíveis eliminações de questões pelo critério do ponto bisserial.
A alternativa correta A, representada no gráfico pela curva em azul, foi escolhida em maiores proporções pelos alunos com desempenho melhor nesta parte da prova.
Já as alternativas incorretas, também denominadas distratores, foram selecionadas principalmente por aqueles com notas mais baixas.
A proporção de alunos que selecionou a resposta correta A aumenta gradativamente, chegando a atingir 100 \% para aqueles que acertaram pelo menos 24 questões, enquanto a proporção dos que escolheram alternativas incorretas decai, a partir da primeira nota não nula ( um pouco mais tarde para a alternativa D ), como função da nota nesta parte da prova.
Os gráficos relativos às demais questões do Componente de Conhecimento Específico constam do Anexo I. 3.3 ANÁLISE DAS QUESTÕES DISCURSIVAS 3.3.1 Componente de Formação Geral As análises dos resultados de desempenho dos estudantes de Medicina nas duas questões discursivas relativas à Formação Geral encontram -se na Tabela 3.8 e no Gráfico 3.15.
Na Tabela 3.8 observa -se que as notas médias foram mais baixas nesse conjunto de questões do que no das objetivas.
Os estudantes de todo o Brasil obtiveram, em Formação Geral, média 60,1 nas questões objetivas e 50,0 nas questões discursivas.
A mediana também confirma o pior desempenho dos alunos de todo o Brasil nas questões discursivas do Componente de Formação Geral.
Enquanto essa estatística foi de 55,5 para questões discursivas, para as questões objetivas a mediana foi de 62,5.
Pode -se, também, notar um aumento do desvio padrão de 19,4, nas questões objetivas do Componente de Formação Geral dos alunos de todo o Brasil, para 21,8 nas questões discursivas do mesmo componente.
Como já comentado, a mediana de todo o Brasil, neste componente, foi 55,5, sendo a menor mediana encontrada na região Norte ( 50,0 ), enquanto a maior mediana foi encontrada na região Centro-Oeste ( 59,5 ).
A maior nota máxima foi obtida, nas questões discursivas do Componente de Formação Geral, por pelo menos um aluno da região Norte ( 97,0 ), enquanto a menor nota máxima foi obtida na região Nordeste ( 92,5 ).
Em todas as regiões a nota mínima foi zero, sem exceção.
O Gráfico 3.15 representa a distribuição das notas nas questões discursivas no Componente de Formação Geral.
A moda desta distribuição ocorre no intervalo ( 50 ; 60 ].
Devido à grande quantidade de notas zero e à alta frequência de alunos que deixaram este tipo de questão em branco o intervalo [ 0 ; 10 ] se apresenta como uma segunda moda.
A distribuição possui assimetria à esquerda, coeficiente de assimetria -1,08.
Em todas as regiões o coeficiente de assimetria também é negativo, o maior em módulo ( -1,30 ), na região Centro-Oeste, e o menor na região Norte, -0,69.
A avaliação das questões discursivas de Formação Geral considerou em separado o conteúdo ( peso 0,8 na nota ) e o desempenho em língua portuguesa ( peso 0,2 ).
Na sequência, os resultados verificados para cada uma das questões discursivas de Formação Geral com relação ao conteúdo serão apresentados, estabelecendo -se relações com os temas abordados em cada uma delas.
Os comentários da Banca de docentes corretores a respeito do observado na correção das respostas dos estudantes, suas impressões e conclusões serão apresentados junto à análise de cada questão.
Em seguida será feita uma análise do desempenho de Língua Portuguesa.
Os comentários da Banca de docentes corretores serão apresentados para o conjunto de questões.
Cumpre esclarecer que, tendo em vista que as questões discursivas de Formação Geral são padronizadas, ou seja, constam de todas as provas, os comentários da Banca são os mesmos para todas as carreiras acadêmicas, sendo direcionados a todos os estudantes que participaram do ENADE/2013.
A seguir, serão analisados os desempenhos da Área de Medicina nas duas questões discursivas de Formação Geral do ENADE/2013, comparando os resultados obtidos com comentários para cada questão. 3.3.1.1 Análise de Conteúdo Questão Discursiva 1 do Componente de Formação Geral Os dados de Medicina, obtidos a partir das respostas à questão 1, encontram -se na Tabela 3.9 e no Gráfico 3.16.
Nessa questão – de melhor desempenho dentre as duas de Formação Geral – os alunos de todo Brasil tiveram média, 51,4.
A maior média para a questão 1 foi obtida na região Centro-Oeste ( 56,7 ), e a menor, na região Norte ( 47,9 ).
Quanto à variabilidade das notas, o desvio padrão de todo o Brasil foi 25,6.
O menor desvio padrão foi obtido na região Centro-Oeste ( 23,4 ) e o maior desvio padrão foi obtido na região Norte ( 28,9 ).
A mediana da questão discursiva 1 foi 55,0 para as regiões Sudeste e Sul, mais alta para as regiões Nordeste e Centro-Oeste ( 60,0 ) e mais baixa para a região Norte ( 50,0 ).
Para o Brasil como um todo foi também 55,0.
As notas máximas ( 100,0 ) foram as mesmas para todas as regiões do Brasil, sem exceção.
Além disso, a nota mínima foi zero em todas as regiões do país, sem exceção.
O Gráfico 3.16 mostra a distribuição das notas na questão discursiva 1 do Componente de Formação Geral.
Observa -se que a maior frequência corresponde aos alunos que obtiveram nota no intervalo ( 60 ; 70 ], correspondendo à moda da distribuição.
Destacam -se, também, os alunos que deixaram esta questão em branco e o intervalo ( 40 ; 50 ], que apresentam máximos locais.
A distribuição das notas possui assimetria negativa ( -0,72 ), assim como em todas as regiões do Brasil ( entre -0,87 para o Centro-Oeste e -0,37 para o Norte ). 3.3.1.2 Comentários sobre a correção de Conteúdo das respostas à Questão Discursiva 1 O enunciado restou claro, elaborado com linguagem técnica, estando inserido nas Diretrizes gerais, com nível médio de profundidade, já que abordava tema vinculado à Área de formação da maior parte dos examinandos.
A divisão das perguntas também ficou clara indicando, passo a passo, o caminho de resposta a seguir, com base no texto.
Verificou -se dificuldade quanto à análise dos programas de governo existentes que contribuem para melhoria do saneamento básico.
Nesse ponto, os alunos referiram, com muita frequência, a ações relacionadas à Saúde, tais como : médico da família, ações educativas, ações ambientais, melhoria nas condições de habitação ( sem necessariamente vincular a programas de governo federal ou estaduais ou municipais ).
Os programas citados foram : Programa Saúde da Família ( PSF ), Programa de Aceleração do Crescimento ( PAC ) e Minha Casa Minha Vida e, ainda, o SUS.
Apesar de uma minoria de examinandos ter indicado o PAC do Governo Federal, a maior parte deles apresentou alguma solução para a questão do saneamento básico.
Os estudantes indicaram, com facilidade, a questão da necessidade do tratamento da água como item essencial do saneamento básico.
Muitos concluintes, através da sua experiência local, citaram situações em que as crianças caminhavam de pés no chão perto de rotas de esgoto a céu aberto.
Também criticaram a inércia governamental em não combater as invasões ilícitas, em locais inadequados, sem estrutura de água potável e também de saneamento básico.
Diante das observações preliminares pode ser extraída uma classificação das respostas obtidas que ficou estabelecida em fracas, médias e boas.
As notas consideradas fracas ( notas de 0 a 3,0 ) foram conferidas àqueles que não conseguiram relacionar o saneamento básico com a importância do tratamento de água, os projetos para resolver o problema e nem indicaram soluções para a questão.
As respostas que tiveram média avaliação ( notas de 3,5 a 7,0 ) tocaram, de algum modo, nos pontos indicados, mesmo que de forma insuficiente.
Já as respostas consideradas boas ( notas de 7,5 a 10,0 ) preencheram os requisitos do padrão de resposta, não logrando nota máxima, na maior parte dos casos, por ausência de indicação de programa governamental, mas sugerindo soluções técnicas para o problema.
As respostas mais comuns foram as consideradas medianas e que indicaram, com frequência, que o problema do saneamento básico estaria nas classes de renda mais desfavorecida e nas regiões consideradas como periferias das grandes cidades.
Também foram frequentes respostas que citaram a necessidade de tratamento de água, acompanhada da necessidade de maiores investimentos nessa área. 3.3.1.3 Análise de Conteúdo da Questão Discursiva 2 do Componente de Formação Geral A Tabela 3.10 mostra que o desempenho dos estudantes na questão 2 ( média 46,4 ) foi inferior ao obtido na questão de número 1 ( média 51,4 ).
A região Centro-Oeste foi aquela onde a média, nessa questão, foi maior ( 51,5 ), e a de menor média foi a região Norte ( 39,9 ).
Quanto à variabilidade das notas, o desvio padrão de todo o Brasil foi 26,5, superior ao obtido na questão de número 1 ( 25,6 ).
O maior desvio nessa questão foi obtido na região Nordeste ( 27,9 ), enquanto o menor foi obtido na região Centro-Oeste ( 24,6 ).
A mediana foi 50,0 em quase todas as regiões, exceto nas regiões Norte ( 45,0 ) e Centro-Oeste ( 55,0 ).
As notas máximas ( 100,0 ) foram as mesmas em todas as regiões do Brasil, assim como as notas mínimas ( 0,0 ), sem exceção.
O Gráfico 3.17 mostra a distribuição das notas na questão discursiva 2 do Componente de Formação Geral.
Neste gráfico, observa -se que a distribuição das notas ignorando -se aquelas questões deixadas em branco ( um máximo local ), é bimodal, com a moda nos intervalos ( 40 ; 50 ] e ( 60 ; 70 ].
O coeficiente de assimetria das notas nesta questão é negativo ( -0,46 ), indicando que a distribuição é assimétrica à esquerda.
O mesmo ocorre para todas as regiões da Brasil ( entre -0,64 na região Centro-Oeste e -0,19 na região Norte ). 3.3.1.4 Comentários sobre a correção de Conteúdo das respostas à Questão Discursiva 2 O tema abordado no enunciado, apesar da complexidade, é extremamente atual.
Essa atualidade facilitou a elaboração de alguma resposta, permitindo que todos os que se dispuseram a responder pudessem apresentar, pelo menos, alguns dos elementos integrantes do padrão de resposta.
A correção, seguindo o padrão de resposta, avaliou a capacidade de elaborar um texto dissertativo que abordasse o tema sob três pontos de vista : o papel da tecnologia digital na atualidade ; a garantia dos direitos do cidadão e do Estado diante do avanço tecnológico ; e o problema da segurança digital para o indivíduo, para o Estado e para as organizações ( empresas e/ou instituições nacionais ou internacionais ).
Dada a amplitude de respostas possíveis, os alunos tiveram desempenho satisfatório.
Identificamos que os concluintes, basicamente, trilharam dois caminhos : a ) a análise da influência das inovações tecnológicas, notadamente através das redes sociais, numa perspectiva essencialmente privada, daí defluindo as questões da privacidade e da segurança ; b ) uma perspectiva de segurança nacional-estatal com criticas à atuação americana de invasão de privacidade dos dados governamentais, e à inercia do governo brasileiro em reagir ao ataque, considerado, pela maioria esmagadora dos concluintes, como indevido e antiético.
Poucos foram os que uniram as duas facetas da questão : o público e o privado.
Na maior parte das provas, o concluinte indicou a ausência de segurança e de privacidade decorrentes das próprias atividades individuais.
Muitos respondentes citaram que, de forma voluntária, os próprios cidadãos enfraquecem a proteção à privacidade, pois divulgam sua vida pessoal, aspectos de seu cotidiano, por meio de imagens e do fornecimento de dados pessoais.
No aspecto segurança, foram inúmeras as indicações de acessos indevidos a dados pessoais com a perpetração de fraudes de todos os tipos, notadamente as bancárias e comerciais, com a aquisição de bens à revelia dos titulares dos dados.
Também no tema segurança, muitos indicaram o aspecto positivo decorrente da vigilância generalizada, não somente através da rede de computadores, mas também pela quantidade de câmeras existentes, que contribuem para a identificação de criminosos procurados ou dos que realizam crimes pela internet.
Na perspectiva pública, a esmagadora maioria indicou repulsa pelo comportamento americano de invasão dos dados do governo brasileiro, reclamando reação do mesmo.
Alguns poucos justificaram a intervenção americana pela “ natural ” disputa de poder pelos Estados imperialistas. 3.3.1.5 Análise de Língua Portuguesa das Questões Discursivas do Componente de Formação Geral Os dados de Medicina, obtidos a partir das respostas às questões discursivas do Componente de Formação Geral no que tange à Língua Portuguesa, encontram -se na Tabela 3.11 e no Gráfico 3.18.
Nesse aspecto, os alunos de todo Brasil tiveram média, 54,2.
O Gráfico 3.18 mostra a distribuição das notas de Língua Portuguesa do Componente de Formação Geral.
Observa -se que a maior frequência corresponde aos alunos que obtiveram nota no intervalo ( 60 ; 70 ], correspondendo à moda da distribuição.
As questões em branco constituem uma segunda moda.
A distribuição das notas possui assimetria negativa ( -1,33 ), assim como em todas as regiões do Brasil ( entre -1,63 para o Centro-Oeste e -1,04 para o Norte ).
A maior média com respeito à Língua Portuguesa foi obtida na região Centro-Oeste ( 58,6 ), e a menor, na região Norte ( 51,9 ).
Quanto à variabilidade das notas, o desvio padrão de todo o Brasil foi 22,3.
O menor desvio padrão foi obtido na região Centro-Oeste ( 18,9 ) e o maior desvio padrão foi obtido na região Nordeste ( 25,1 ).
A mediana da nota de Língua Portuguesa foi 60,0 para as regiões Sudeste e Sul, mais alta para as regiões Nordeste e Centro-Oeste ( 62,5 ) e mais baixa para a região Norte ( 57,5 ).
Para o Brasil como um todo foi também 60,0.
A maior nota máxima foi obtida por pelo menos um aluno na região Nordeste ( 95,0 ), enquanto a menor nota máxima foi obtida na região Norte ( 90,0 ).
Além disso, a nota mínima foi zero em todas as regiões do país, sem exceção.
O Gráfico 3.18 mostra a distribuição das notas de Língua Portuguesa do Componente de Formação Geral.
Observa -se que a maior frequência corresponde aos alunos que obtiveram nota no intervalo ( 60 ; 70 ], correspondendo à moda da distribuição.
As questões em branco constituem uma segunda moda.
A distribuição das notas possui assimetria negativa ( -1,33 ), assim como em todas as regiões do Brasil ( entre -1,63 para o Centro-Oeste e -1,04 para o Norte ). 3.3.1.6 Comentários sobre a correção das respostas de Formação Geral com respeito à Língua Portuguesa Um aspecto interessante a ser destacado é a diferença de desempenho entre a questão 1 e a questão 2, principalmente nos aspectos textuais – os textos produzidos para a questão 1 apresentaram resultados melhores do que os produzidos para a questão 2.
Uma hipótese que provavelmente explique essa discrepância é o fato de o tema “ saneamento básico ” ser mais relevante para profissionais de saúde do que o tema “ segurança e privacidade ”, constituindo tema trabalhado no currículo escolar.
Prova disso é que quase todos os textos da questão 1 citam termos específicos da área de saúde, como “ contaminação feco-oral ”, “ higiene corporal ”, “ ingestão de alimentos ou água contaminada por vetores biológicos ”.
Os resultados da avaliação correspondem aos seguintes aspectos observados em cada competência : a ) aspectos ortográficos : o desempenho dos participantes ficou concentrado, majoritariamente, nos níveis 3 ( até 3 desvios ) e 2 ( de 4 a 7 desvios ), devido ao grande índice de desvios de acentuação, contrabalançada pelo pequeno número de desvios em grafia.
Foi reduzido o número de textos com ausência de desvios ( enquadrados no nível 4 ).
Alguns textos, com pior desempenho, foram enquadrados no nível 1 ( de 8 a 12 desvios ), enquanto o nível zero foi atribuído a pouquíssimos casos, já que ele revela falta absoluta de domínio das convenções ortográficas.
Observou -se, portanto, que existe uma diferença muito grande de desempenho nos dois aspectos analisados : baixíssimo índice de desvios ortográficos e grande índice de desvios de acentuação.
Em vários casos, ocorre ausência completa de acentuação gráfica.
Os resultados revelam, portanto, que a tendência dominante entre os universitários brasileiros é a eliminação da acentuação gráfica, talvez motivada pela ausência de esclarecimento dos meios de comunicação, das autoridades e das escolas sobre as decisões do Acordo Ortográfico da Língua Portuguesa de 1990.
Os casos mais sistemáticos de ausência de acentuação são : • palavras proparoxítonas ( “ politicas ”, “ publicas ”, “ juridicas ” ) ; • palavras paroxítonas terminadas em ditongo crescente ( “ varios ”, “ Brasilia ”, “ familia ” ) ; • palavras oxítonas ( “ ninguem ”, “ esta ”, “ ate ” ).
Quanto ao domínio das convenções relativas à grafia das palavras, são pontuais os problemas encontrados ( “ infecciozas ”, “ obtos ”, “ hipedemia ”, “ saniamento ”, “ cituação ”, “ abitantes ” ).
Os únicos desvios recorrentes foram : a ausência de segmentação nos casos de verbo + pronome oblíquo “ se ”, combinada, eventualmente, à alteração da grafia do pronome ( “ percebese ”, “ percebece ”, “ preferece ” ) ; a grafia da forma verbal de pretérito imperfeito do subjuntivo como se fosse uma ocorrência de verbo + pronome oblíquo ( “ distribui -se ” por “ distribuísse ” ) Vale observar, também, que, ao contrário do que se esperava, não apareceram abreviaturas próprias do “ internetês ”, ou seja, dos hábitos de comunicação escrita adquiridos pelo uso de redes sociais e e-mails. b ) aspectos textuais : o desempenho dos participantes ficou concentrado, majoritariamente, nos níveis 3 e 2, devido à grande ocorrência de problemas de estruturação textual.
Foi muito baixo o número de textos enquadrados no nível 4 por não apresentar qualquer problema estrutural.
Alguns textos com pior desempenho foram enquadrados no nível 1, enquanto o nível zero foi atribuído a pouquíssimos casos, já que ele revela a existência de um texto sem articulação e com comprometimento do sentido.
Observou -se que a grande maioria dos participantes não distribuiu as ideias em parágrafos, talvez devido ao pequeno número de linhas disponibilizadas para a resposta da questão ou, quem sabe, pela suposição de que não seria necessária essa divisão por não se tratar de um texto no modelo de uma redação dissertativo-argumentativa, como solicitado nos vestibulares.
Essa competência é a que se revela como a mais problemática entre os participantes, porque são muitos os problemas observados, desvios acumulados durante toda a formação do estudante e que não se resolvem com um estudo autodidata, como acontece com regras ortográficas ou morfossintáticas : sequência justaposta de ideias sem encaixamentos sintáticos ; redução drástica de estruturas subordinadas, ao lado do aumento na frequência de estruturas coordenadas e absolutas ; redução no uso de conectores para expressar relações lógicas essenciais à construção do texto, substituídas pela exigência de inferência por parte do interlocutor para suprir a sua ausência ; emprego equivocado de operadores que não estabelecem relações lógicas coerentes entre ideias do texto ; emprego inadequado do pronome relativo ( com omissão da preposição ou a utilização de pronome inadequado, como “ onde ” ) ; repetição exaustiva de termos sem a utilização de procedimentos mais sofisticados de substituição ( hiperonímias, hiponímias, nominalizações, expressões metafóricas ) ; frases fragmentadas que comprometem a estrutura lógico-gramatical ; frases formadas apenas por oração subordinada, sem oração principal.
Um importante aspecto a destacar é o baixíssimo desempenho de uma parte dos participantes em relação à estrutura formal do texto produzido, o que é extremamente preocupante ao se levar em conta que são graduandos em fase final de formação.
São frequentes os casos de desvios de estruturação frasal, com uso inadequado ou ausência de conectivos entre parágrafos e entre frases.
Em uma parte dos textos, falta um mínimo de textualidade e de domínio do registro padrão da língua.
Na verdade, observam -se relações linguísticas quase agramaticais, como as estabelecidas pela sequência de gerúndios sem o apoio de um ponto de partida para a organização das informações gramaticais e semânticas.
Quanto à utilização dos mecanismos de referenciação, dois aspectos devem ser destacados : a ocorrência, em uma boa parte dos textos, de repetições de palavras ou expressões sem a utilização de termos sinônimos ou pronomes, como seria adequado ; e a alternância entre os pronomes demonstrativos este / esse, e suas formas flexionadas, sem a observância de um critério de utilização.
Esse último aspecto tem se generalizado na modalidade oral da língua e, cada vez mais, tem se tornado uma constante na modalidade escrita de jornais e revistas : predominância da forma “ esse ” na modalidade oral e da forma “ este ” na modalidade escrita, provavelmente por um processo de hipercorreção. c ) aspectos morfossintáticos e vocabulares : o desempenho dos participantes ficou concentrado, majoritariamente, nos níveis 3 e 2, devido à grande ocorrência de problemas de regência, concordância e uso adequado dos sinais de pontuação.
O nível 4 foi atribuído a um número muito reduzido de textos.
Alguns textos com pior desempenho foram enquadrados no nível 1, enquanto o nível zero foi atribuído a pouquíssimos casos, já que ele revela a existência de um texto sem articulação e com comprometimento do sentido.
Os resultados são muito transparentes em relação aos aspectos mais problemáticos no desempenho dos participantes.
O desvio mais frequente é a falta do sinal indicativo da crase, o que se configura como um problema de regência, verbal ou nominal, conforme o caso – isso revela que o usuário não tem consciência de que, sob a forma do termo “ a ” existe a presença de uma preposição “ a ”, exigida pela regência do termo anterior.
Embora em outros exames, como o Enem, a falta de crase seja apenada em acentuação, nesta correção amostral esse desvio foi considerado no âmbito dos aspectos morfossintáticos.
Outro problema relacionado à regência verbal e à nominal, encontrado frequentemente nas questões, foi a ausência de preposição antes de pronome relativo, processo generalizado na modalidade oral da língua, em situações de registro informal.
Apesar da possibilidade de que essa alteração de regência se generalize no padrão escrito da Língua Portuguesa, como já está ocorrendo até em textos jornalísticos, o não emprego da preposição foi apenado neste processo de avaliação.
Outro desvio muito frequente diz respeito aos processos de concordância verbal e de concordância nominal.
Quanto à concordância de número, observou -se ausência de marca ( com sujeito anteposto ou posposto ) ou uso indevido ( uso inadequado da marca de plural comandado pelo núcleo plural da locução adjetiva, apesar de o substantivo que funciona como núcleo do sintagma nominal estar no singular ).
Uma ocorrência que se destacou foi a ausência de acento circunflexo na forma plural do presente do indicativo do verbo “ ter ”, que foi considerada como um desvio na concordância verbal e não na acentuação gráfica.
Quanto à concordância de gênero, vários casos foram observados, normalmente no âmbito de sintagmas nominais longos, em que o adjetivo está afastado do substantivo.
Quanto à questão da colocação pronominal, foram poucos os casos observados.
Apesar de serem aspectos relacionados à oralidade, concluiu -se que, no registro escrito formal, a maioria dos participantes já incorporou regras como a não introdução da frase por um pronome oblíquo e a próclise na presença de um termo atrator.
Não se adotou, entretanto, o padrão excessivamente formal descrito pelas gramáticas normativas em relação à posição do pronome oblíquo em locuções verbais, já que esse uso está muito distante da prática cotidiana, até em textos mais formais.
Quanto aos aspectos vocabulares, alguns tipos de inadequação foram observados : expressões da oralidade apareceram em algumas respostas, mas sem maior relevância do ponto de vista quantitativo ; seleção vocabular incompatível com o contexto, gerando situações de falta de inteligibilidade ; falta de domínio de vocabulário mais abstrato e de maior complexidade, essencial ao desenvolvimento do texto de base dissertativa.
Quanto à utilização dos sinais de pontuação, observou -se uma grande precariedade nos textos analisados.
Não foi apenada a ausência de vírgula para destacar locuções ou adjuntos adverbiais de pequena extensão deslocados de posição na frase, por ser um uso opcional.
São os seguintes os tipos de problemas encontrados : a ) vírgula : utilização de vírgula para separar o sujeito e o predicado ; ocorrência de apenas uma das vírgulas para separar uma palavra, uma expressão ou uma oração encaixada ; uso de vírgula no lugar do ponto para separar ideias que constituem períodos distintos ; ausência de vírgula para separar enumerações ; ausência de vírgula para separar oração adjetiva explicativa ou utilização inadequada para separar oração adjetiva restritiva ; b ) ponto e vírgula : utilização do ponto e vírgula no lugar de vírgula ; c ) ponto final : ausência de ponto final para separar períodos.
Várias marcas de oralidade foram identificadas, embora não em alta frequência : o uso do pronome relativo “ onde ” como relativo universal, falta de artigo definido antes de substantivo, repetição de palavras por falta de vocabulário, reduções como “ tá ”, “ pra ”, “ pro ”, “ prum ”, expressões informais como “ estão nem aí ”, “ tem que ”, “ levando com a barriga ”, “ não tá nem aí ”, “ a coisa fica séria ”, “ e aí haja dor de cabeça ”. 3.3.2 Componente de Conhecimento Específico Na parte da prova relativa às questões discursivas no Componente de Conhecimento Específico ( Tabela 3.12 ), observa -se que a média foi mais baixa do que nas questões discursivas do Componente de Formação Geral.
Enquanto no Componente de Formação Geral a média para estudantes de Medicina de todo o Brasil foi 50,0, na Parte de Conhecimento Específico a média foi 22,9.
A maior média deste componente foi obtida pelos estudantes da região Centro-Oeste ( 26,9 ), e a menor, pelos da região Norte ( 16,4 ).
Quanto à variabilidade das notas, o desvio padrão de todo o Brasil foi 15,8.
O maior desvio padrão foi encontrado na região Centro-Oeste ( 16,7 ), e o menor, na região Norte ( 12,8 ).
A maior nota máxima, nas questões discursivas do Componente de Conhecimento Específico, foi obtida na região Sudeste ( 89,2 ), enquanto a menor nota máxima foi encontrada na região Norte ( 72,5 ).
Além disso, a nota mínima ( 0,0 ) foi obtida por alunos de todas as regiões do Brasil sem exceção.
A mediana do Brasil como um todo foi 20,0, também encontrada nas regiões Nordeste e Sudeste.
As demais medianas foram : 14,2 na região Norte, 25,0 na região Sul e 24,2 na região Centro-Oeste.
O Gráfico 3.19 representa a distribuição das notas nas questões discursivas do Componente de Conhecimento Específico.
Essa distribuição tem moda no intervalo de ( 10 ; 20 ].
Nota -se que há uma tendência decrescente a partir deste intervalo até o último intervalo.
O coeficiente de assimetria é positivo para todo o Brasil ( 0,74 ) e para todas as regiões ( entre 0,40 para a região Sul e 1,03 para a região Norte ), indicando uma assimetria à direita. 3.3.2.1 Análise da Questão Discursiva 3 do Componente de Conhecimento Específico Na questão 3, cujos resultados aferidos encontram -se descritos na Tabela 3.13, a média dos estudantes de todo o Brasil foi 29,8.
O desempenho médio dos estudantes de todo o Brasil nesta questão foi superior ao desempenho nas questões 4 e 5.
A menor média nessa questão foi obtida pelos alunos da região Norte ( 24,7 ), enquanto a maior média foi obtida na região Centro-Oeste ( 33,6 ).
Quanto à variabilidade das notas, o desvio padrão de todo o Brasil foi 21,0.
O maior desvio padrão foi obtido na região Centro-Oeste ( 21,6 ), enquanto o menor foi obtido na região Nordeste ( 19,8 ).
A maior nota máxima foi alcançada por pelo menos um aluno tanto na região Norte quanto na região Sudeste ( 97,5 ), enquanto a menor nota máxima foi obtida nas regiões Sul e Centro-Oeste ( 90,0 ).
A mediana foi 30,0 para as regiões Nordeste, Sudeste e Sul, mais alta para a região Centro-Oeste ( 32,5 ) e mais baixa para a região Norte ( 20,0 ).
Para o Brasil como um todo foi também 30,0.
A nota mínima do Brasil como um todo foi zero, o mesmo valor para todas as regiões, sem exceção.
O Gráfico 3.20 mostra a distribuição das notas na questão discursiva 3 do Componente de Conhecimento Específico.
Essa distribuição é unimodal com moda no intervalo ( 20 ; 30 ].
O coeficiente de assimetria é positivo para todo o Brasil ( 0,27 ) e para todas as regiões ( entre 0,09 na região Sul e 0,64 na região Norte ), indicando uma assimetria à direita. 3.3.2.2 Comentários sobre as respostas à Questão Discursiva 3 O conteúdo e o nível de profundidade exigido na Questão 3 atendiam ao que preconizam as Diretrizes Curriculares.
O enunciado era claro quanto à descrição do caso e redigido em linguagem adequada.
No entanto, no item ‘ c ’, houve uso de linguagem inadequada na expressão “ demonstrando a dose ”, o que não prejudicou o entendimento do que era solicitado.
Apesar de ser possível haver interpretações diferentes da esperada, a principal hipótese diagnóstica era claramente identificável a partir da descrição do caso.
O padrão de respostas garantia a consideração da diversidade de respostas nos itens ‘ a ’ e ‘ c ’.
As situações mais comuns foram erros de dose e menção correta da via sem menção do local de aplicação.
Não identificamos, exceto quanto ao erro de dose já mencionado, outras situações em que o padrão não previu respostas dadas pelos estudantes.
Alguns tipos de respostas foram encontrados com maior frequência :  A pergunta solicitava dados da anamnese e a resposta se referia a exame físico.
Ou seja, um número significativo de alunos confunde anamnese com exame físico.  A medicação estava correta, mas a via de aplicação errada e não havia explicação sobre o mecanismo de ação.  O aluno, talvez por recurso mnemônico para ' decorar ', apresenta os itens da anamnese, mas erra o diagnóstico de gravidade do caso.
Também houve os que acertaram o tratamento, mas erraram o mecanismo de ação.
A impressão geral do desempenho foi a de que os alunos tiveram mais sucesso na identificação do quadro clínico do que em relação à terapia.
O grande número de notas baixas diante de um quadro clínico de fácil identificação e cuja ocorrência não é rara, é um resultado preocupante. 3.3.2.3 Análise da Questão Discursiva 4 do Componente de Conhecimento Específico A Tabela 3.14 contém as informações relativas à questão 4 do conjunto de questões do Componente de Conhecimento Específico.
O desempenho dos estudantes de todo o Brasil nesta questão foi inferior ao desempenho nas questões de número 3 e 5.
A média geral do Brasil foi 10,9, sendo a menor média registrada na região Norte ( 4,7 ) e a maior na região Centro-Oeste ( 13,9 ).
Quanto à variabilidade das notas, o desvio padrão de todo o Brasil foi 19,9.
O maior desvio padrão foi obtido na região Centro-Oeste ( 22,3 ), enquanto o menor foi obtido na região Norte ( 12,6 ).
A nota máxima ( 100,0 ) foi atingida em todas as regiões, sem exceção.
A mediana e a nota mínima foram zero para todas as regiões do Brasil, sem exceção.
Ou seja, pelo menos metade dos alunos tirou nota zero nesta questão.
O Gráfico 3.21 representa a distribuição de notas na questão discursiva 4 do Componente de Conhecimento Específico.
A moda desta distribuição ocorre nos alunos que tentaram resolver esta questão e receberam nota zero.
Destaca -se, também, a alta frequência de alunos que deixaram esta questão em branco ( cerca de 27 \% ).
O coeficiente de assimetria é positivo para todo o Brasil ( 2,13 ) e para todas as regiões ( entre 1,74 para a região Sul e 3,77 para a região Norte ), indicando uma assimetria acentuada à direita. 3.3.2.4 Comentários sobre as respostas à Questão Discursiva 4 O conteúdo da Questão 4 era adequado em relação às diretrizes curriculares, com o nível de profundidade é médio/fácil.
O enunciado, quanto à linguagem era razoavelmente claro, mas não em relação ao conteúdo.
A questão apresentava sérios problemas em relação à possibilidade de diferentes interpretações.
Por exemplo, o que foi solicitado no item ‘ a ’ independia totalmente da tabela, no entanto, o caput solicitava que os alunos respondessem considerando as informações da tabela, o que acabou por confundir.
O item ‘ b ’ permitia a interpretação de que o que se desejava era apenas uma explicação geral da interpretação do resultado de qualquer RR, do ponto de vista da significância estatística.
Não foi solicitado, explicitamente, que o aluno identificasse os valores do RR que tinham significância estatística na tabela e desse sua interpretação a esses valores.
Infelizmente, problemas quanto ao enunciado não podem ser corrigidos depois que a prova já foi aplicada.
No entanto, esses problemas se refletiram nas respostas redigidas pelos alunos.
Apesar das dificuldades diagnosticadas em relação ao enunciado da questão e ao padrão de resposta, o que ficou claro, a partir dos resultados da correção, é que o conteúdo avaliado não está sendo ministrado adequadamente e de forma plena a todos os alunos dos cursos de Medicina.
A impressão geral do desempenho foi de que um número muito pequeno de alunos domina o conteúdo, e a grande maioria não tem conhecimento mínimo sobre o assunto.
Foi grande o número de notas baixas e de notas zero enquanto foi bem pequeno o de notas altas.
Algumas situações foram encontradas com maior frequência :  Confusão entre RR e taxas de mortalidade ; letalidade ; incidência etc.  RR definido como : “ o risco comparado com o risco da população em geral ”  Tentativa não bem sucedida de ‘ descobrir ’ a resposta a partir da tabela.
Em geral, podemos dizer que a maioria das provas revela que esse assunto é desconhecido.
Algumas provas apontaram para a possibilidade de RR ser igual a zero, outras incluem a necessidade de p valor para responder aos quesitos ; etc. 3.3.2.5 Análise da Questão Discursiva 5 do Componente de Conhecimento Específico A Tabela 3.15 contém as informações relativas à questão 5 do conjunto do Componente de Conhecimento Específico.
O desempenho dos estudantes nessa questão foi inferior ao da questão 3 e superior ao da questão 4.
A nota média dos estudantes de todo o Brasil foi 28,0.
A maior média foi registrada nas regiões Sul e Centro-Oeste ( 33,3 ), enquanto a menor média foi registrada na região Norte ( 19,9 ).
Quanto à variabilidade das notas, o desvio padrão dos alunos do Brasil, como um todo, foi 21,7.
Enquanto o maior desvio foi encontrado na região Nordeste ( 23,2 ), o menor foi encontrado na região Norte ( 17,9 ).
A nota máxima ( 100,0 ) foi alcançada em todas as regiões do Brasil, sem exceção.
A mediana foi 25,0 para a região Nordeste, mais alta ( 30,0 ) para as regiões Sul e Centro-Oeste e mais baixa para as regiões Norte ( 15,0 ) e Sudeste ( 20,0 ).
Para o Brasil como um todo foi também 25,0.
A nota mínima foi zero para todas as regiões do Brasil, sem exceção.
O Gráfico 3.22 representa a distribuição das notas da questão discursiva 5 do Componente de Conhecimento Específico.
A moda desta distribuição ocorre no intervalo ( 10 ; 20 ].
Destacam -se, também, os alunos que deixaram esta questão em branco aparecendo como uma moda secundária.
O coeficiente de assimetria para todos os alunos é 0,69 e para todas as regiões se mantém positivo, variando entre 0,32 na região Sul e 1,00 na região Norte. 3.3.2.6 Comentários sobre as respostas à Questão Discursiva 5 O enunciado da Questão 5 era extenso, porém claro e com linguagem adequada.
O conteúdo avaliado está entre os previstos nas Diretrizes Curriculares e o nível de profundidade da questão era médio/grande.
A descrição do quadro clínico era bastante clara e não permitia interpretações diversas da esperada.
O padrão de resposta tem linguagem adequada e foi adequado à correção da questão, e o nível de profundidade exigido era compatível.
Por sua abrangência, não houve resposta dada pelos estudantes, que pudesse estar correta, e não estivesse prevista.
A impressão geral do desempenho é de que os alunos tiveram mais sucesso na identificação dos valores normais dos parâmetros solicitados.
A identificação da característica da dor e do problema do “ pé diabético ” foi feita pela minoria.
A isso corresponde o erro frequente da terapia da dor e das recomendações dos “ cuidados gerais ”.
No entanto, é curioso ressaltar a impressão de que o erro em relação às recomendações voltadas para o “ pé diabético ” tenha sido muito mais frequente do que o erro em relação à terapia da dor.
Tal como na Questão 3, o grande número de notas baixas, diante de um quadro clínico de fácil identificação, foi muito frequente, e este é um resultado preocupante.
Dentre as respostas redigidas algumas chamaram atenção por serem mais frequentes : • Item ‘ a ’ – avaliação da PA e triglicerídeos correta com cálculo errado do LDL ; alteração do tratamento errada ; item ‘ b ’ parcialmente correto – carbamazepina e... paracetamol ; idem item ‘ c ’ – auto exame dos pés e... dieta. • item ‘ b ’ em branco ou totalmente errado.
Em geral, os erros mais frequentes se referem à terapia da dor e à prevenção do “ pé diabético ”.
Houve um número grande de notas baixas e um pequeno número de notas altas.

% matched lemmas: 
\end{textsample}
